\section{Psihološki mehanizmi i posledice po korisnike}

Korišćenjem znanja o ljudskim psihološkim mehanizmima, tehnička optimizacija pažnje stvara konkretne korisničke navike: konstantno proveravanje, skrolovanje, čekanje reakcije i sve teže prekidanje interakcije. Korisnik nije pasivan objekat, ali nije ni potpuno autonoman akter koji pred sobom ima neutralan ekran i slobodno odlučuje koliko će ostati. Nalazi se u prostoru koji je dizajniran tako da prepoznaje njegove reakcije, slabosti, navike i emocionalne okidače, a zatim ih koristi za optimizaciju zadržavanja pažnje.

Problem nije u tome što platforme poznaju ljudsko ponašanje. Problem je u tome što se to znanje često koristi unutar poslovnog modela u kome se korisnikova pažnja ne štiti, već eksploatiše. Kada je cilj duže zadržavanje, češći povratak i intenzivnija reakcija, psihologija prestaje da bude samo nauka o čoveku i postaje alat za dizajniranje okruženja iz kog se čovek sve teže sam izvlači.

\subsection{Promenljive nagrade i uklanjanje tačke prekida}

\paragraph{Neizvesnost nagrade.}
Jedan od osnovnih mehanizama na kojima počiva adiktivnost digitalnih platformi jeste neizvesnost nagrade~\cite{montag2019addictive}. Korisnik ne zna šta ga čeka kada otvori aplikaciju, osveži tok sadržaja ili nastavi da skroluje. Možda ga čeka poruka, zanimljiv video, smešna objava, nova reakcija na njegov sadržaj ili informacija koja mu se čini važnom. Možda ga ne čeka ništa posebno. Korisnik ne dolazi samo zbog sadržaja koji je već video, već zbog mogućnosti da sledeći sadržaj bude baš onaj koji će ga dotaći.

Ovakva logika promenljive nagrade posebno je moćna jer ne traži stalno zadovoljstvo. Naprotiv, sistem ne mora uvek da pruži korisniku nešto prijatno da bi ga zadržao. Dovoljno je da povremeno pruži dovoljno jak signal da se nastavak isplati. Tok sadržaja zato ne mora biti ravnomerno zanimljiv; on može biti mešavina dosade, iznenađenja, smeha, zavisti, šoka, straha, tuge i konflikta. Korisnik nastavlja ne zato što je svaki sadržaj vredan, već zato što sledeći možda jeste.

\paragraph{Uklanjanje tačke prekida.}
Drugi važan mehanizam jeste uklanjanje tačke prekida. Tradicionalni mediji uglavnom imaju jasne krajeve: emisija se završi, novine se pročitaju. Digitalne platforme taj kraj sistematski razgrađuju. Beskonačno skrolovanje (engl. \textit{infinite scroll}), automatsko pokretanje sadržaja (engl. \textit{autoplay}) i beskonačni tok sadržaja uklanjaju trenutak u kome bi korisnik prirodno zastao i doneo odluku da prestane~\cite{montag2019addictive}. Umesto da korisnik bira sledeći sadržaj, sistem mu ga već servira. Umesto da se pojavi praznina u kojoj bi mogao da zatvori aplikaciju, pojavljuje se novi stimulans.

Zato beskonačni tok sadržaja nije samo udobna funkcionalnost. On je arhitektura produženog ostanka. Kada nema kraja, korisnik mora sam da proizvede prekid, a upravo je to najteže u prostoru koji je podešen da svaki pokušaj prekida dočeka novom mogućnošću, novom notifikacijom, novim sadržajem i novim obećanjem da se još malo ostane.

\subsection{Socijalna validacija, strah od propuštanja (FOMO) i emocionalna promenljivost toka sadržaja}

\paragraph{Socijalna validacija.}
Posebno snažan sloj psihološke eksploatacije nastaje onda kada platforma ne optimizuje samo odnos korisnika prema sadržaju, već i njegov odnos prema drugim ljudima. Lajkovi, komentari, deljenja, pregledi i pratioci nisu samo brojevi. Oni postaju javni znaci vrednosti, prihvaćenosti i društvene vidljivosti~\cite{wadsley2022reward}. Korisnik ne proverava aplikaciju samo da bi video sadržaj, već i da bi proverio kako je on sam prošao u očima drugih.

Zato socijalna validacija ima dvojaku funkciju. Sa jedne strane, ona korisniku daje osećaj povezanosti, priznanja i prisustva. Sa druge strane, ona ga vezuje za sistem u kome je njegov društveni značaj stalno merljiv, uporediv i nesiguran. Objavljena slika, komentar ili video ne završavaju se trenutkom objavljivanja. Oni nastavljaju da žive kroz broj reakcija, a korisnik nastavlja da proverava da li je dobio dovoljno pažnje. Platforma tako ne meri samo šta korisnik gleda, već i kako korisnik sebe vidi kroz reakcije drugih.

\paragraph{Strah od propuštanja (engl. \textit{fear of missing out}, FOMO).}
Strah od propuštanja (FOMO) dodatno pojačava ovu dinamiku. Korisnik se ne vraća samo zato što želi nešto konkretno, već zato što ne želi da propusti nešto što se možda dešava bez njega. Platforma tada ne prodaje samo sadržaj, već osećaj da je odsustvo rizično. Biti van aplikacije znači potencijalno zakasniti, ne znati, ne učestvovati, ne biti viđen.

\paragraph{Emocionalna promenljivost toka sadržaja.}
Emocionalna promenljivost toka sadržaja dodatno održava pažnju. Korisniku se ne nudi nužno uvek isti ton, jer bi predvidljivost brzo postala dosadna. Upravo smena prijatnog, smešnog, tužnog, šokantnog, konfliktnog i intimnog sadržaja održava korisnika budnim. Srećan kratki video-snimak (engl. \textit{reel}) može biti praćen tužnim, nežan sadržaj agresivnim, informativan senzacionalističkim. U ovom okruženju pažnja se ne zadržava samo interesovanjem, već i emocionalnom destabilizacijom, a takav tok sadržaja nije neutralna lista objava, već psihološki ritam podešen tako da korisnik ne potone u dosadu dovoljno dugo da bi aplikaciju zatvorio.

\subsection{Od navike ka digitalnoj adikciji}

\paragraph{Gubitak kontrole.}
Granica između navike i digitalne adikcije nije uvek jasna, ali se može prepoznati po gubitku kontrole. Nije problem samo u tome što korisnik često koristi platformu, već u tome što sve teže prestaje i kada zna da mu korišćenje ne prija. Nastavlja da skroluje i kada je umoran, proverava aplikaciju i kada nema konkretan razlog, vraća se i kada ga sadržaj očigledno iscrpljuje.

\paragraph{Odgovornost korisnika i odgovornost sistema.}
Potrebno je ponovo naglasiti da digitalna adikcija ne nastaje samo iz slabosti pojedinca. Takvo objašnjenje je previše zgodno za platforme, jer svu odgovornost prebacuje na korisnika. Naravno da čovek ima odgovornost za sopstvene navike, ali ona ne može biti analizirana van konteksta sistema koji je tehnički, ekonomski i psihološki podešen da prekid učini težim, a povratak verovatnijim. Problem nije u samom dizajnu, već u cilju kojem dizajn služi. Kada cilj nije dobrobit korisnika, nego maksimizacija njegove angažovanosti, granica između pomoći i eksploatacije postaje namerno mutna.

\paragraph{Dizajn kao sistem.}
Zbog toga je naivno posmatrati aplikacije isključivo kao zabavu ili praktične alate. One su deo infrastrukture pažnje. Na njihovom dizajnu rade ljudi koji razumeju ponašanje, navike, emocije, nagrade i prekide pažnje. To znanje nije neutralno kada se primeni u sistemu čiji profit raste sa vremenom koje korisnik provodi unutra. Ako kompanija zna da notifikacija vraća korisnika, da beskonačni tok sadržaja produžava sesiju, da socijalna validacija izaziva proveravanje, a da promenljiva nagrada održava iščekivanje, onda se više ne može govoriti o slučajnosti. To nije greška sistema. To je sistem.

\paragraph{Posledice po pažnju i odnos prema sebi.}
Posledice se ne završavaju na izgubljenom vremenu. Dugotrajno izlaganje ovakvom okruženju menja način na koji korisnik raspoređuje pažnju, doživljava dosadu, podnosi tišinu i gradi odnos prema sebi i drugima. Ako je svaki prazan trenutak odmah popunjen sadržajem, dosada prestaje da bude prostor odmora, mišljenja ili kreativnosti i postaje signal za otvaranje aplikacije. Ako je društvena vrednost stalno merena reakcijama, samopouzdanje postaje zavisno od spoljašnje metrike. Ako je tok sadržaja emocionalno nestabilan, korisnik se navikava na stanje stalne pobuđenosti.

U tome leži najdublja posledica psihološke eksploatacije: platforme ne uzimaju korisniku samo vreme, već preoblikuju njegov odnos prema sopstvenoj pažnji. Čovek sve teže ostaje sam sa sobom, sve teže trpi odsustvo stimulacije i sve lakše prihvata da svaka praznina mora biti popunjena sadržajem. U poslovnom smislu, to je idealan korisnik: prisutan, reaktivan, merljiv i povratljiv. U ljudskom smislu, to je iscrpljen čovek čija pažnja postaje korporativno vlasništvo.